\documentclass[11pt,a4paper]{article}
\usepackage[utf8]{inputenc}
\usepackage[T1]{fontenc}
\usepackage{geometry}
\geometry{margin=1in}
\usepackage{hyperref}
\usepackage{enumitem}
\usepackage{booktabs}
\usepackage{longtable}
\usepackage{xcolor}
\usepackage{titlesec}
\usepackage{amsmath,amssymb}

\titleformat{\section}{\Large\bfseries\color{blue!70!black}}{}{0em}{}
\titleformat{\subsection}{\large\bfseries\color{blue!50!black}}{}{0em}{}
\titleformat{\subsubsection}{\normalsize\bfseries\color{blue!30!black}}{}{0em}{}

\title{\textbf{Replication Plan}\\[0.5em]
\large Common-Mode Voltage Reduction of Single-Phase Quasi-Z-Source Inverter-Based Photovoltaic System\\[0.3em]
\normalsize OSTI ID: 1606674}
\author{Auto-generated Replication Blueprint}
\date{\today}

\begin{document}
\maketitle
\tableofcontents
\newpage

%---------------------------------------------------------------------
\section{Paper Overview}
%---------------------------------------------------------------------
This paper proposes and analyzes modulation strategies for reducing common-mode voltage (CMV) in a single-phase quasi-Z-source inverter (qZSI) used in transformerless photovoltaic (PV) grid-tied systems. The high CMV in conventional modulation causes undesirable ground leakage current through the PV panel parasitic capacitance. The authors derive the CMV expressions for different switching states, propose modified PWM schemes that avoid high-CMV states, and validate through circuit simulation showing reduced CMV and leakage current while maintaining output power quality.

%---------------------------------------------------------------------
\section{Detailed Workflow}
%---------------------------------------------------------------------

\subsection{Phase 1: Circuit Topology Modeling}
\begin{enumerate}[label=\textbf{1.\arabic*}]
  \item \textbf{Model the quasi-Z-source inverter.}
    \begin{itemize}
      \item Define the qZSI topology: two inductors $L_1$, $L_2$; two capacitors $C_1$, $C_2$; diode $D$; H-bridge (4 switches $S_1$--$S_4$).
      \item Model the shoot-through and non-shoot-through states.
      \item Define circuit parameters from the paper (inductance, capacitance, switching frequency, DC input voltage).
    \end{itemize}
  \item \textbf{Model the PV parasitic capacitance} ($C_{\text{PV}}$) and grid connection.
  \item \textbf{Define switching states} and their corresponding CMV expressions.
\end{enumerate}

\subsection{Phase 2: Modulation Strategy Implementation}
\begin{enumerate}[label=\textbf{2.\arabic*}]
  \item \textbf{Implement conventional unipolar and bipolar PWM.}
  \item \textbf{Implement the proposed CMV-reduced PWM scheme.}
    \begin{itemize}
      \item Identify and eliminate high-CMV switching states.
      \item Generate gate signals for $S_1$--$S_4$ according to the modified modulation logic.
      \item Incorporate shoot-through intervals for voltage boost.
    \end{itemize}
  \item \textbf{Implement comparison schemes} (e.g., H5 modulation, HERIC-type) if discussed.
\end{enumerate}

\subsection{Phase 3: Circuit Simulation}
\begin{enumerate}[label=\textbf{3.\arabic*}]
  \item \textbf{Build the circuit model} in MATLAB/Simulink (Simscape Power Systems) or PLECS or LTspice.
  \item \textbf{Run time-domain simulations} for each modulation strategy.
    \begin{itemize}
      \item Simulation time: several grid cycles (e.g., 0.2\,s at 60\,Hz).
      \item Record: output voltage/current, CMV waveform, leakage current through $C_{\text{PV}}$.
    \end{itemize}
  \item \textbf{Compute metrics.}
    \begin{itemize}
      \item CMV peak and RMS values.
      \item Leakage current peak and RMS.
      \item Output voltage THD (total harmonic distortion).
      \item Efficiency and power quality metrics.
    \end{itemize}
\end{enumerate}

\subsection{Phase 4: Results Reproduction}
\begin{enumerate}[label=\textbf{4.\arabic*}]
  \item Reproduce CMV waveform comparison plots.
  \item Reproduce leakage current waveforms.
  \item Reproduce output voltage/current THD comparison.
  \item Compare tabulated numerical results against paper.
\end{enumerate}

%---------------------------------------------------------------------
\section{Datasets}
%---------------------------------------------------------------------
No external datasets; all data is generated from circuit simulations with parameters specified in the paper.

\begin{longtable}{p{4.5cm} p{3.5cm} p{6cm}}
\toprule
\textbf{Dataset / Source} & \textbf{Type} & \textbf{Location / Notes} \\
\midrule
\endfirsthead
\toprule
\textbf{Dataset / Source} & \textbf{Type} & \textbf{Location / Notes} \\
\midrule
\endhead
Circuit parameters & Specified in paper & Table of component values ($L$, $C$, $f_{\text{sw}}$, $V_{\text{DC}}$, etc.) \\
\midrule
Switching state tables & Specified in paper & CMV expressions for each state combination \\
\midrule
Grid parameters & Standard & 120\,V/240\,V, 60\,Hz (or as specified) \\
\bottomrule
\end{longtable}

%---------------------------------------------------------------------
\section{Models, Tools, and Software}
%---------------------------------------------------------------------

\subsection{Simulation Platforms}
\begin{longtable}{p{4.5cm} p{2.5cm} p{6.5cm}}
\toprule
\textbf{Tool / Library} & \textbf{Version} & \textbf{Purpose / URL} \\
\midrule
\endfirsthead
\toprule
\textbf{Tool / Library} & \textbf{Version} & \textbf{Purpose / URL} \\
\midrule
\endhead
MATLAB/Simulink & R2021b+ & Circuit simulation with Simscape Electrical \newline \url{https://www.mathworks.com/} \\
\midrule
PLECS & latest & Alternative power electronics simulator \newline \url{https://www.plexim.com/plecs} \\
\midrule
LTspice & latest & Free SPICE simulator (alternative) \newline \url{https://www.analog.com/en/resources/design-tools-and-calculators/ltspice-simulator.html} \\
\midrule
Python + PySpice & latest & Open-source SPICE interface \newline \url{https://pyspice.fabrice-salvaire.fr/} \\
\midrule
NumPy / SciPy & $\geq 1.21$ & Numerical analysis of waveforms, FFT for THD \\
\midrule
Matplotlib & $\geq 3.4$ & Waveform plotting \\
\bottomrule
\end{longtable}

\subsection{Hardware Requirements}
\begin{itemize}
  \item Standard workstation; circuit simulations are not computationally demanding.
\end{itemize}

\end{document}
